% The only source of this talk. `make TALK=indeterminate-beams` builds slides,
% notes, script, handout, transparencies and article from it. See README.md.
%   - Text outside frames appears only in the article.
%   - Every frame carries a \note{...} for the presenter.
%   - Numbers come from data/numbers.tex (generated by derive.py), never typed.
\documentclass[
  beameroptions={aspectratio=169,ignorenonframetext,11pt},
  articleoptions={11pt,a4paper},
]{beamerswitch}
\usepackage{yjtalk}
% Generated by derive.py. Do not edit.
\newcommand{\ibL}{6}
\newcommand{\ibW}{10}
\newcommand{\ibEI}{16\,000}
\newcommand{\ibWL}{60}
\newcommand{\ibCantMA}{180}
\newcommand{\ibDeltaW}{101.25}
\newcommand{\ibRB}{22.5}
\newcommand{\ibRA}{37.5}
\newcommand{\ibMA}{45}
\newcommand{\ibXmax}{3.75}
\newcommand{\ibMmax}{25.3125}
\newcommand{\ibMmaxShort}{25.3}
\newcommand{\ibXcf}{1.5}
\newcommand{\ibXmaxFromB}{2.25}
\newcommand{\ibLone}{4}
\newcommand{\ibLtwo}{6}
\newcommand{\ibLsum}{10}
\newcommand{\ibWtotal}{100}
\newcommand{\ibRhsOne}{160}
\newcommand{\ibRhsTwo}{540}
\newcommand{\ibRhsSum}{700}
\newcommand{\ibLhsCoef}{20}
\newcommand{\ibMB}{35}
\newcommand{\ibTRA}{11.25}
\newcommand{\ibTRB}{64.58}
\newcommand{\ibTRC}{24.17}
\newcommand{\ibTRBl}{28.75}
\newcommand{\ibTRBr}{35.83}
\newcommand{\ibTxOne}{1.125}
\newcommand{\ibTMOne}{6.33}
\newcommand{\ibTxTwo}{2.42}
\newcommand{\ibTMTwo}{29.2}
\newcommand{\ibDFBA}{0.6}
\newcommand{\ibDFBC}{0.4}
\newcommand{\ibFEMBA}{20}
\newcommand{\ibFEMBC}{-45}
\newcommand{\ibUnbal}{-25}
\newcommand{\ibUnbalAbs}{25}
\newcommand{\ibDistBA}{15}
\newcommand{\ibDistBC}{10}
\newcommand{\ibMBA}{35}
\newcommand{\ibMBC}{-35}
\newcommand{\ibEqualMB}{45}
\newcommand{\ibTxTwoFromA}{7.58}
\newcommand{\ibTZeroOne}{2.25}
\newcommand{\ibTHogOne}{1.75}


\title{Statically indeterminate beams}
\subtitle{When equilibrium runs out, compatibility supplies the missing equations}
\author{Teo Yu Jie}
\institute{A recap, with two worked examples}
\date{30 September 2026}

\pgfplotstableread[col sep=comma]{data/propped.csv}\propped
\pgfplotstableread[col sep=comma]{data/twospan.csv}\twospan

% A beam drawn from (0,0) to (#1,0) in TikZ units; supports and loads are added per frame.
\tikzset{
  beam/.style={line width=2.2pt, yjForeground},
  udl/.style={-Stealth, yjWarm, thin},
  react/.style={-Stealth, yjAccent, thick},
}
\newcommand{\pinsupport}[1]{%
  \draw[yjSecondary, fill=yjSurface] (#1,0) -- ++(-0.18,-0.3) -- ++(0.36,0) -- cycle;}
\newcommand{\fixedsupport}[1]{%
  \draw[yjSecondary, line width=1pt] (#1,-0.45) -- (#1,0.45);
  \foreach \yy in {-0.4,-0.2,0,0.2,0.4} {\draw[yjSecondary] (#1,\yy) -- ++(-0.15,-0.12);}}
\newcommand{\udlarrows}[2]{% from x=#1 to x=#2
  \draw[yjWarm] (#1,0.55) -- (#2,0.55);
  \foreach \xx in {#1,...,#2} {\draw[udl] (\xx,0.55) -- (\xx,0.05);}}

\begin{document}

\begin{frame}[plain]
  \titlepage
  \narration{Some beams have more supports than equilibrium can pay for. This
    recap shows where the extra equations come from, works two examples, and
    checks every number twice.}
  \note{Open with the promise: by the end you can count the unknowns, write
    the missing equation, and check the answer three ways. About twelve
    minutes. Every number on the slides is generated by derive.py and checked
    against the beamdiag solver.}
\end{frame}

\mode<article>{\maketitle}

\section{Counting}

A planar beam has three independent equilibrium equations: forces along the
beam, forces across it, and moments. When the supports provide more reaction
components than that, statics alone cannot find them. The surplus is the
degree of static indeterminacy. For a single beam, with $r$ reaction
components, it is $r-3$ (Hibbeler, \emph{Structural Analysis}, counts
determinacy the same way). When there is no axial load, the horizontal
equation and the horizontal reactions both drop out, so the count is unchanged
for a propped cantilever or a continuous beam, but a beam fixed at both ends
then has two redundants for bending, not three.

\begin{frame}{Equilibrium gives three equations; extra supports add unknowns it cannot find}
  \centering\small
  \begin{tabular}{lccc}
    \toprule
    Beam & Reactions $r$ & Equations & Degree $r-3$ \\
    \midrule
    Simply supported (pin + roller) & 3 & 3 & 0 \\ \pause
    Propped cantilever (fixed + roller) & 4 & 3 & 1 \\ \pause
    Two-span continuous (pin + 2 rollers) & 4 & 3 & 1 \\ \pause
    Fixed at both ends & 6 & 3 & 3 \\
    \bottomrule
  \end{tabular}
  \par\bigskip
  \uncover<4->{Each redundant needs one \yjkey{compatibility} equation: a
    deflection or slope the supports force to be zero.}
  \narration{A beam in a plane gives three equilibrium equations. A simply
    supported beam has three reactions, so it is determinate. A propped
    cantilever has four, one too many. So does a beam continuous over three
    supports. A beam fixed at both ends has six. Each surplus reaction needs one
    extra equation, and it comes from compatibility.}
  \note{Stress that the degree is a count, not a difficulty. Point at the
    fixed--fixed row: with no axial load its horizontal reactions vanish, so
    for bending it has two redundants. Transition: where does the extra
    equation come from?}
\end{frame}

\section{The force method}

The force method, also called the flexibility or compatibility method, makes
the beam determinate by removing one support reaction, the redundant. The
released beam, the primary structure, now deflects where the real support
allows no deflection. Superposing the deflection due to the loads and the
deflection due to the redundant, and setting their sum to the real value,
gives one equation per redundant. The method needs only the deflection formulas
of determinate beams (tabulated in, for example, Gere and Goodno,
\emph{Mechanics of Materials}).

\begin{frame}{Release one reaction, then make the released beam fit its support again}
  \begin{enumerate}
    \item Count the redundants and pick one: here the prop force $R_B$.
    \item<2-> Remove it. The primary structure is a cantilever, which is determinate.
    \item<3-> Find the tip deflection under the load, $\delta_{B,w}$, and per unit $R_B$, $f_{BB}$.
    \item<4-> Compatibility: the real prop does not move, so $\delta_{B,w} - R_B\,f_{BB} = 0$.
    \item<5-> Equilibrium now gives every other reaction.
  \end{enumerate}
  \narration{Pick one reaction as the redundant, here the force in the prop.
    Remove it, and the beam becomes a cantilever, which statics can solve.
    Find how far the tip falls under the load, and how far it rises per unit
    prop force. The real prop does not move, so the two must cancel. Then
    equilibrium gives everything else.}
  \note{Five steps, one per click. The choice of redundant is free; choosing
    the moment at the wall instead gives a simply supported primary beam and
    the same answer. Likely question: does the sign matter? Only consistently:
    here down is the load's deflection, up is the prop's. Transition: do it
    with numbers.}
\end{frame}

\begin{frame}{The prop carries three eighths of the load: $R_B = 3wL/8$}
  \begin{columns}[c]
    \column{0.42\textwidth}
      \centering
      \begin{tikzpicture}[x=0.8cm, y=0.8cm]
        \udlarrows{0}{6}
        \draw[beam] (0,0) -- (6,0);
        \fixedsupport{0}
        \pinsupport{6}
        \node[font=\footnotesize, text=yjWarm, above] at (3,0.55) {$w = \ibW$ kN/m};
        \draw[react] (6,-1.4) -- (6,-0.35) node[pos=0, below, font=\footnotesize, text=yjAccentText] {$R_B$};
        \draw[yjSecondary, |-|] (0,-0.8) -- (6,-0.8) node[midway, below, font=\footnotesize] {$L = \ibL$ m};
        \node[font=\footnotesize, text=yjSecondary] at (0,-1.3) {$A$};
      \end{tikzpicture}
    \column{0.55\textwidth}
      \begin{align*}
        \delta_{B,w} &= \frac{wL^4}{8EI} = \ibDeltaW~\text{mm} \\
        \uncover<2->{f_{BB} &= \frac{L^3}{3EI} \\}
        \uncover<3->{R_B &= \frac{\delta_{B,w}}{f_{BB}} = \frac{3wL}{8} = \yjkey{\ibRB~\text{kN}}}
      \end{align*}
      \uncover<4->{\small Then $R_A = wL - R_B = \ibRA$~kN and
        $M_A = wL^2/8 = \ibMA$~kN\,m, hogging.\par}
  \end{columns}
  \yjsource{$EI=\ibEI$ kN\,m$^2$ ($E=200$ GPa, $I=8\times10^{-5}$ m$^4$); $EI$ cancels in $R_B$.}
  \narration{A cantilever under a uniform load sags w L to the fourth over
    eight E I at the tip: here, \ibDeltaW\ millimetres. A unit tip force lifts
    it L cubed over three E I. Setting the two equal gives the prop force:
    three eighths of the load, \ibRB\ kilonewtons. Equilibrium then gives
    \ibRA\ kilonewtons and a hogging moment of \ibMA\ kilonewton metres at the
    wall.}
  \note{Point at how EI cancels: the reactions of a uniform beam do not
    depend on its stiffness, only the deflections do. The two deflection
    formulas are standard cantilever results. Transition: what did the prop
    buy us?}
\end{frame}

\begin{frame}{The prop cuts the peak moment from \ibCantMA\ to \ibMA~kN\,m}
  \centering
  \begin{tikzpicture}
    \begin{axis}[yj, width=12cm, height=4.8cm, xmin=0, xmax=6, ymin=-190, ymax=40,
        xlabel={distance from the wall $x$ (m)}, ylabel={bending moment $M$ (kN\,m)}]
      \addplot[yjBorder, line width=1pt, dashed] table[x=x, y=cant] {\propped};
      \node[anchor=west, font=\footnotesize, text=yjSecondary] at (axis cs:0.3,-165) {no prop: $-wL^2/2 = -\ibCantMA$ at the wall};
      \only<2->{
        \addplot[yjAccent, line width=1.2pt] table[x=x, y=prop] {\propped};
        \fill[yjAccent] (axis cs:\ibXmax,\ibMmax) circle[radius=2.5pt];
        \node[anchor=south, font=\footnotesize, text=yjAccentText] at (axis cs:\ibXmax,\ibMmax) {$+\ibMmax$ at $x=\ibXmax$};
        \node[anchor=west, font=\footnotesize, text=yjAccentText] at (axis cs:0.1,-45) {$-\ibMA$};
      }
      \only<3->{
        \fill[yjWarm] (axis cs:\ibXcf,0) circle[radius=2.5pt];
        \node[anchor=north west, font=\footnotesize, text=yjWarm] at (axis cs:\ibXcf,0) {$M=0$ at $x=\ibXcf=L/4$};
      }
      \draw[yjSecondary, very thin] (axis cs:0,0) -- (axis cs:6,0);
    \end{axis}
  \end{tikzpicture}
  \yjsource{Sagging positive: $M(x) = -\ibMA + \ibRA\,x - \ibW\,x^2/2$ kN\,m.}
  \narration{Without the prop, the wall carries a moment of \ibCantMA\ kilonewton
    metres. With it, the wall moment falls to \ibMA, and the largest sagging
    moment is \ibMmaxShort, where the shear is zero, \ibXmax\ metres from the
    wall. The moment changes sign a quarter of the way along.}
  \note{Point at the three marks: the wall moment, the sagging peak where shear
    is zero, and the point of contraflexure at $L/4$, where reinforcement or a
    splice would go. The peak is $9wL^2/128$. Transition: what about a beam
    that runs over several supports?}
\end{frame}

\section{Continuous beams}

For a beam continuous over several supports, the convenient redundants are the
bending moments over the interior supports. Cutting the beam into simply
supported spans at those supports releases the continuity of slope; restoring
it at each interior support gives Clapeyron's three-moment equation (1857),
derived for example in Megson, \emph{Structural and Stress Analysis}. The
form below assumes a uniform $EI$, supports at one level and a uniform load
$w_i$ on span $i$.

\begin{frame}{The three-moment equation makes the slope continuous over each support}
  \[ M_{A} L_1 + 2 M_B (L_1 + L_2) + M_C L_2 = -\frac{w_1 L_1^3}{4} - \frac{w_2 L_2^3}{4} \]
  \begin{itemize}
    \item<2-> Redundants: the support moments. Each span becomes simply supported.
    \item<3-> Each right-hand term is $6EI$ times a simply supported end slope, $wL^3/24EI$.
    \item<4-> Pinned ends: $M_A = M_C = 0$. One interior support, one equation.
  \end{itemize}
  \yjsource{Uniform $EI$, supports at one level, sagging moments positive (Clapeyron, 1857).}
  \narration{For a beam over several supports, take the moments over the
    supports as the redundants. Each span becomes simply supported, and the
    slopes on either side of a support must match. That gives the three moment
    equation. At pinned ends the moment is zero, so one interior support means
    one equation.}
  \note{Stress the physical meaning: it is a slope-matching condition, one per
    interior support. Point at the right-hand side: $wL^3/24EI$ is the end
    slope of a simply supported span, times $6EI$. Transition: apply it to two
    unequal spans.}
\end{frame}

\begin{frame}{Two spans of \ibLone\ and \ibLtwo\ m: the support moment is \ibMB~kN\,m, hogging}
  \begin{columns}[c]
    \column{0.44\textwidth}
      \centering
      \begin{tikzpicture}[x=0.5cm, y=0.8cm]
        \udlarrows{0}{10}
        \draw[beam] (0,0) -- (10,0);
        \pinsupport{0} \pinsupport{4} \pinsupport{10}
        \node[font=\footnotesize, text=yjWarm, above] at (5,0.55) {$w = \ibW$ kN/m};
        \foreach \xx/\nn in {0/A, 4/B, 10/C} {\node[font=\footnotesize, text=yjSecondary, below] at (\xx,-0.35) {$\nn$};}
        \draw[yjSecondary, |-|] (0,-1.1) -- (4,-1.1) node[midway, below, font=\footnotesize] {\ibLone\ m};
        \draw[yjSecondary, |-|] (4,-1.1) -- (10,-1.1) node[midway, below, font=\footnotesize] {\ibLtwo\ m};
      \end{tikzpicture}
    \column{0.54\textwidth}
      \begin{align*}
        2M_B(\ibLone + \ibLtwo) &= -\tfrac{\ibW\cdot \ibLone^3}{4} - \tfrac{\ibW\cdot \ibLtwo^3}{4} \\
        \uncover<2->{\ibLhsCoef\,M_B &= -\ibRhsOne - \ibRhsTwo = -\ibRhsSum \\}
        \uncover<3->{M_B &= \yjkey{-\ibMB~\text{kN\,m}}}
      \end{align*}
      \uncover<4->{\small Span by span: $R_A = \ibTRA$, $R_B = \ibTRBl + \ibTRBr = \ibTRB$,
        $R_C = \ibTRC$ kN.\par}
  \end{columns}
  \narration{Two spans, of \ibLone\ and \ibLtwo\ metres, carry \ibW\ kilonewtons
    per metre. The three moment equation reads \ibLhsCoef\ times M B equals
    minus \ibRhsSum. So the moment over the middle support is
    \ibMB\ kilonewton metres, hogging. Each span, taken on its own with that end
    moment, gives the reactions.}
  \note{Walk the arithmetic: $\ibW\cdot\ibLone^3/4 = \ibRhsOne$ and
    $\ibW\cdot\ibLtwo^3/4 = \ibRhsTwo$. Reactions come from each span as a free
    body with $M_B$ applied: $R_A = wL_1/2 + M_B/L_1$. Point out that $R_B$ is
    well over half the \ibWtotal\ kN load. Transition: the diagram.}
\end{frame}

\begin{frame}{The long span sags most; the support moment is the largest}
  \centering
  \begin{tikzpicture}
    \begin{axis}[yj, width=12cm, height=4.5cm, xmin=0, xmax=10, ymin=-42, ymax=38,
        xlabel={distance from $A$ (m)}, ylabel={bending moment $M$ (kN\,m)}]
      \addplot[yjAccent, line width=1.2pt] table[x=x, y=M] {\twospan};
      \draw[yjSecondary, very thin] (axis cs:0,0) -- (axis cs:10,0);
      \fill[yjAccent] (axis cs:4,-\ibMB) circle[radius=2.5pt];
      \node[anchor=west, font=\footnotesize, text=yjAccentText] at (axis cs:4.15,-\ibMB) {$M_B = -\ibMB$};
      \node[anchor=south, font=\footnotesize, text=yjSecondary] at (axis cs:\ibTxOne,\ibTMOne) {$+\ibTMOne$};
      \node[anchor=south, font=\footnotesize, text=yjSecondary] at (axis cs:\ibTxTwoFromA,\ibTMTwo) {$+\ibTMTwo$};
    \end{axis}
  \end{tikzpicture}
  \yjsource{Sagging positive. Peaks at zero shear: \ibTxOne\ m from $A$, \ibTxTwo\ m from $C$.}
  \narration{The moment over the support, \ibMB\ kilonewton metres, is the
    largest in the beam. The long span still sags by up to \ibTMTwo, and the
    short span, held down by the support moment, sags only \ibTMOne.}
  \note{Point at the support peak, then at the two sagging peaks. The short
    span hogs over its last \ibTHogOne\ m, from $x=\ibTZeroOne$ m to $B$: the
    long span's load pulls it up over the support. Transition: how do we know $M_B$ is right?}
\end{frame}

\section{Checks}

Moment distribution (Cross, 1930) reaches the same support moment by a
different route. Lock the interior joint against rotation, so each span
starts with its fixed-end moments; with the far end pinned these are
$wL^2/8$. Then release the joint: the unbalanced moment is shared between the
members in proportion to their stiffness, $3EI/L$ for a member whose far end
is pinned. With only one joint and pinned far ends there is nothing to carry
over, so a single balance is exact.

\begin{frame}{Moment distribution reaches the same \ibMB~kN\,m in one balance}
  \centering
  \begin{tabular}{lrr}
    \toprule
    At joint $B$ & member $BA$ & member $BC$ \\
    \midrule
    Stiffness $3EI/L$, share & \ibDFBA & \ibDFBC \\ \pause
    Fixed-end moment $\pm wL^2/8$ & $+\ibFEMBA$ & $\ibFEMBC$ \\ \pause
    Balance the $\ibUnbal$ unbalance & $+\ibDistBA$ & $+\ibDistBC$ \\ \pause
    Final end moment & $\yjkey{+\ibMBA}$ & $\yjkey{\ibMBC}$ \\
    \bottomrule
  \end{tabular}
  \par\bigskip
  \uncover<4->{End moments are clockwise-positive; $\pm\ibMB$ is the \ibMB\ kN\,m
    hogging moment found by the three-moment equation.\par}
  \yjsource{Far ends pinned: modified stiffness $3EI/L$, no carry-over (Cross, 1930).}
  \narration{Moment distribution gets there another way. Lock joint B, and each
    span starts with its fixed end moment. Release the joint, and an
    unbalance of \ibUnbalAbs\ kilonewton metres is shared in proportion to stiffness: \ibDistBA\ to
    the short span and \ibDistBC\ to the long one. The end moments come out
    at \ibMB, the same answer.}
  \note{Stress that this is an independent method, not a rearrangement: it
    uses stiffnesses, not flexibilities. The shares are $(3/4)/(3/4+3/6)$ and
    $(3/6)/(3/4+3/6)$. Transition: four checks worth running every time.}
\end{frame}

\begin{frame}{Four checks catch most mistakes before they reach a drawing}
  \begin{enumerate}
    \item Vertical equilibrium: $\ibTRA + \ibTRB + \ibTRC = \ibWtotal$ kN, and moments about $A$ balance.
    \item<2-> Limit case: equal spans must give $M_B = -wL^2/8$ (\ibEqualMB\ kN\,m for $L=\ibLtwo$ m).
    \item<3-> Compatibility: the deflection comes out zero at every support.
    \item<4-> An independent solver: the site's beam visualizer, by the stiffness method, matches every number here.
  \end{enumerate}
  \yjsource{Checks 3 and 4: beamdiag's exact-arithmetic reference solver, both examples.}
  \narration{Run four checks. The reactions must add up to the load. Equal
    spans must give w L squared over eight. The deflection must be zero at
    every support. And an independent solver must agree: the beam visualizer
    on this site matches every number in this talk.}
  \note{Stress that equilibrium alone cannot confirm an indeterminate answer:
    any wrong $R_B$ still satisfies it once the other reactions follow. The
    compatibility check is the one that proves the redundant. Transition: what
    to remember.}
\end{frame}

\begin{frame}{What to remember}
  \begin{enumerate}
    \item Count first: the degree of indeterminacy is reactions minus three.
    \item<2-> One compatibility equation per redundant: release it, then make the beam fit again.
    \item<3-> For continuous beams, support moments and the three-moment equation are quickest.
    \item<4-> Check with a second method and with compatibility, not only with equilibrium.
  \end{enumerate}
  \vfill
  {\scriptsize\color{yjSecondary}
    Hibbeler, \emph{Structural Analysis}.\quad
    Gere \& Goodno, \emph{Mechanics of Materials}.\quad
    Megson, \emph{Structural and Stress Analysis}.\quad
    Clapeyron (1857), Comptes rendus 45.\quad
    Cross (1930), \emph{Analysis of continuous frames by distributing fixed-end moments}, Proc.\ ASCE 56.\par}
  \narration{Remember four things. Count the unknowns first. Write one
    compatibility equation per redundant. For continuous beams, use the
    support moments. And check with a second method, not only with
    equilibrium.}
  \note{Say the four lines, one per click, then stop. Closing line: statics
    tells you what balances, compatibility tells you what fits, and an
    indeterminate beam needs both. Point people to the beam visualizer to try
    their own spans. Take questions.}
\end{frame}

\end{document}
